% Generated draft: requires mathematical and source review.
\documentclass[11pt]{article}
\usepackage[margin=25mm]{geometry}
\usepackage{fontspec}
\setmainfont{Latin Modern Roman}
\usepackage{amsmath,amssymb,mathtools,unicode-math}
\setmathfont{Latin Modern Math}
\usepackage{enumitem,needspace,xurl,hyperref}
\setlength{\parindent}{0pt}
\setlength{\parskip}{5pt}
\setlength{\emergencystretch}{3em}
\newcommand{\N}{\mathbb{N}}
\newcommand{\Z}{\mathbb{Z}}
\newcommand{\Q}{\mathbb{Q}}
\newcommand{\R}{\mathbb{R}}
\newcommand{\C}{\mathbb{C}}
\newcommand{\F}{\mathbb{F}}
\newcommand{\A}{\mathbb{A}}
\newcommand{\PP}{\mathbb{P}}
\newcommand{\E}{\mathbb{E}}
\newcommand{\eps}{\varepsilon}
\newcommand{\dd}{\,\mathrm{d}}
\begin{document}
\section*{UnsolvedMath ID: 700446}
\textbf{Source problem code:} catalogue entry
\subsection{Definitions and mathematical statement}
Let \(\Omega\) be a set with at least two elements and let \(G\leq\operatorname{Sym}(\Omega)\). The group \(G\) is locally finite if every finitely generated subgroup of \(G\) is finite. Its action is primitive if it is transitive and the only \(G\)-invariant equivalence relations on \(\Omega\) are equality and the universal relation. No finiteness or cardinality bound on \(\Omega\) is imposed.

For this draft, a Sylow \(2\)-subgroup of a locally finite group means a subgroup maximal by inclusion among its \(2\)-subgroups, where a \(2\)-subgroup is one whose elements all have finite order a power of \(2\). This uses the customary infinite-group convention, which the supplied statement does not explicitly define. For a subgroup \(P\), write
\[
Z(P)=\{z\in P:zp=pz\text{ for every }p\in P\}.
\]
An involution is an element of order exactly \(2\), and
\[
\operatorname{Fix}_{\Omega}(t)=\{\omega\in\Omega:t(\omega)=\omega\}.
\]

\textbf{Classification problem (A. N. Fomin, Kourovka Notebook, Problem 7.52).} Describe all locally finite primitive permutation groups \(G\leq\operatorname{Sym}(\Omega)\) satisfying
\[
\forall P\leq G\text{ maximal among }2\text{-subgroups},\quad
\exists t\in Z(P)\quad
\bigl(t\neq 1,\ t^2=1,\ |\operatorname{Fix}_{\Omega}(t)|=1\bigr).
\]
Thus each Sylow \(2\)-subgroup must contain at least one central involution fixing exactly one point; the condition does not require every central involution to do so.

The supplied source reports that D. Holt completely determined the finite case in 1978. This is retained as a source qualification, without independent verification or an identification of the resulting finite groups. The source does not specify the equivalence relation intended for the requested classification.
\subsection{Short English statement}
Describe the locally finite primitive permutation groups for which every Sylow \(2\)-subgroup has a central involution fixing exactly one point. The supplied record attributes the problem to A. N. Fomin and reports that D. Holt determined the finite case in 1978; the remaining classification is not given.
\subsection{Sources}
\begin{itemize}
\item Supplied local UnsolvedMath registry record, permanent ID \(700446\), exact source code \texttt{catalogue entry}, title \emph{Kourovka Notebook Problem 7.52}. The registry labels the problem \texttt{open}; this is supplied metadata, not an independently verified status assessment.
\item \emph{The Kourovka Notebook (Unsolved Problems in Group Theory)}, Problem 7.52, proposed by A. N. Fomin, 7th issue (1980). The supplied background identifies the consulted edition as the 21st issue edition (2026), arXiv:1401.0300v48. The supplied collection description names E. I. Khukhro and V. D. Mazurov as editors.
\item The supplied statement attributes the complete determination of the finite case to D. Holt in 1978. No separate bibliographic reference for that result is supplied.
\end{itemize}
\end{document}
